% Secciones corregidas para reemplazar desde "Tecnologias Utilizadas" hasta
% "Integracion del Sistema" en el capitulo Desarrollo del Sistema.
% El resto del capitulo (introduccion, requerimientos y diseno) puede conservarse.

\subsection{Tecnologías Utilizadas}

La implementación del sistema requirió tecnologías orientadas al desarrollo
web, la persistencia relacional, la extracción de datos y el procesamiento
asíncrono. La selección definitiva quedó conformada por las siguientes
herramientas:

\begin{itemize}
    \item \textit{Frontend}: Vue.js 3, TypeScript, Vite, Pinia y Tailwind CSS
    \cite{vuejs}.
    \item \textit{Backend}: Python 3.11 o superior y FastAPI
    \cite{pythondocs,fastapi}.
    \item \textit{Persistencia}: SQLAlchemy como ORM y Alembic para el control
    versionado del esquema \cite{sqlalchemy}.
    \item \textit{Base de datos de demostración}: SQLite, utilizada para el
    desarrollo local, las migraciones y las pruebas de integración.
    \item \textit{Base de datos objetivo de producción}: PostgreSQL, soportada
    por la configuración y los adaptadores de persistencia, pero pendiente de
    validación integral en un ambiente productivo \cite{postgresdocs}.
    \item \textit{Validación y serialización}: Pydantic \cite{pydantic}.
    \item \textit{Scraping HTTP}: \texttt{aiohttp}, utilizado para consumir de
    forma asíncrona los catálogos y API públicas disponibles.
    \item \textit{Scraping con navegador}: Playwright, utilizado para fuentes
    cuyo contenido requiere ejecución de JavaScript \cite{playwright}.
    \item \textit{Concurrencia}: \texttt{asyncio}, mediante grupos de tareas,
    semáforos, colas y límites temporales \cite{pythondocs}.
    \item \textit{Pruebas y calidad}: pytest, pytest-asyncio, Vitest y Ruff.
    \item \textit{Control de versiones}: Git y GitHub \cite{github}.
\end{itemize}

No se utilizaron Requests, Beautiful Soup ni \texttt{httpx} en los adaptadores
de scraping vigentes. Los catálogos estáticos o públicos se consultan como datos
JSON mediante \texttt{aiohttp}, mientras que La Anónima se procesa mediante
Playwright.

\subsection{Implementación del Backend}

El backend fue desarrollado con FastAPI y organizado como un monolito modular
con principios de Arquitectura Limpia y Arquitectura Hexagonal. Los módulos
\texttt{catalog}, \texttt{supermarkets}, \texttt{prices}, \texttt{basket},
\texttt{geo}, \texttt{decision} e \texttt{ingestion} separan entidades de
dominio, casos de uso, puertos, adaptadores de infraestructura e interfaces HTTP.

La lógica de negocio no depende directamente de FastAPI ni de SQLAlchemy. Los
casos de uso operan mediante puertos y una Unidad de Trabajo que coordina los
repositorios y las transacciones. Esta organización permite reemplazar SQLite
por PostgreSQL sin modificar el dominio ni el modelo de decisión.

La API pública utiliza el prefijo \texttt{/api/v1}. Sus operaciones principales
se resumen en la Tabla~\ref{tab:endpoints_api_reales}.

\begin{table}[H]
\centering
\small
\begin{tabular}{|p{5.6cm}|p{1.5cm}|p{6cm}|}
\hline
\textbf{Endpoint} & \textbf{Método} & \textbf{Función} \\
\hline
\texttt{/health} & GET & Consultar el estado básico del backend. \\
\hline
\texttt{/api/v1/catalog/products} & GET & Buscar productos del catálogo normalizado. \\
\hline
\texttt{/api/v1/catalog/categories} & GET & Consultar categorías disponibles. \\
\hline
\texttt{/api/v1/catalog/brands} & GET & Consultar marcas registradas. \\
\hline
\texttt{/api/v1/locations/cities} & GET & Consultar ciudades y coordenadas de referencia. \\
\hline
\texttt{/api/v1/supermarkets} & GET & Consultar cadenas de supermercados. \\
\hline
\texttt{/api/v1/branches} & GET & Consultar sucursales por ciudad o supermercado. \\
\hline
\texttt{/api/v1/prices/current} & GET & Consultar precios vigentes y métricas de calidad. \\
\hline
\texttt{/api/v1/prices/history} & GET & Consultar el historial de una publicación. \\
\hline
\texttt{/api/v1/prices/compare} & GET & Comparar precios de productos normalizados. \\
\hline
\texttt{/api/v1/basket/validate} & POST & Validar una canasta temporal. \\
\hline
\texttt{/api/v1/decisions/ranking} & POST & Generar el ranking multicriterio. \\
\hline
\texttt{/api/v1/ingestion/sources} & GET/POST & Listar o crear fuentes de extracción. \\
\hline
\texttt{/api/v1/ingestion/sources/refresh-concurrently} & POST & Actualizar varias fuentes concurrentemente. \\
\hline
\end{tabular}
\caption{Principales endpoints implementados en la API v1}
\label{tab:endpoints_api_reales}
\end{table}

Además de estas operaciones, la API de ingestión permite actualizar la
configuración de una fuente, consultar corridas y finalizar una corrida como
exitosa o fallida. Estas funciones proporcionan trazabilidad administrativa y
no forman parte del flujo de compra del usuario final.

\subsection{Implementación del Scraping}

El módulo de scraping utiliza un contrato común que permite incorporar fuentes
sin modificar los casos de uso de ingestión. Cada fuente configurada conserva
el supermercado asociado, la sucursal de destino, la URL base, el adaptador a
utilizar y su estado de activación.

La demostración integra cuatro adaptadores reales:

\begin{itemize}
    \item \textit{Carrefour}: consulta su infraestructura VTEX mediante
    \texttt{aiohttp}. El adaptador crea una sesión de compra, asocia el código
    postal configurado y registra si la ubicación fue confirmada.
    \item \textit{Jumbo}: consulta el catálogo público mediante solicitudes
    HTTP asíncronas.
    \item \textit{La Coope}: consulta su API pública y filtra resultados que no
    contienen los términos significativos de búsqueda.
    \item \textit{La Anónima}: utiliza Playwright porque el catálogo depende de
    contenido dinámico ejecutado en el navegador.
\end{itemize}

Todos los adaptadores devuelven el mismo contrato de producto extraído, que
incluye código externo, identificador comercial, nombre, marca, precio,
presentación, URL, imagen, fuente y ciudad declarada. Los resultados se validan
y se almacenan primero en una tabla de staging antes de modificar el catálogo o
el historial de precios.

El alcance actual corresponde a un piloto para Comodoro Rivadavia y un conjunto
reducido de productos. Los catálogos públicos de Jumbo y La Coope no permiten
confirmar por sí solos una sucursal física. Por este motivo, sus precios deben
interpretarse como datos del piloto asociado y no como cobertura nacional ni
como evidencia de disponibilidad en cualquier local de la cadena.

\subsection{Implementación de la Concurrencia y Alcance del Paralelismo}

Las tareas de extracción están limitadas principalmente por operaciones de
entrada/salida, como la latencia de red y el tiempo de respuesta de los sitios
externos. En consecuencia, se adoptó concurrencia asíncrona en lugar de
paralelismo de CPU.

El caso de uso de actualización concurrente implementa los siguientes
mecanismos:

\begin{itemize}
    \item \texttt{asyncio.TaskGroup} para administrar las tareas por fuente;
    \item \texttt{asyncio.Semaphore} para limitar la concurrencia máxima;
    \item \texttt{asyncio.Queue} para entregar resultados de extracción;
    \item límites temporales por fuente mediante \texttt{asyncio.timeout};
    \item aislamiento de fallos, de modo que el error de una fuente no cancela
    los resultados válidos de las restantes;
    \item persistencia ETL controlada y serializada después de cada extracción.
\end{itemize}

Los adaptadores HTTP también reutilizan una sesión y limitan las consultas
simultáneas. Para La Anónima se implementó un pool acotado de páginas: una única
instancia de Chromium mantiene varios contextos y páginas reutilizables, cada
una asignada de forma exclusiva mientras procesa una consulta. Por lo tanto, no
se crea un navegador independiente por cada worker.

En Windows, Playwright puede ejecutarse en un hilo dedicado con un bucle Proactor
cuando el servidor utiliza un bucle incompatible con la creación de subprocesos.
Este mecanismo resuelve una restricción de integración y no constituye
paralelismo de cálculo intensivo.

El benchmark final realizado con Carrefour, La Coope y La Anónima registró tres
repeticiones completas por modalidad. El promedio secuencial fue de
aproximadamente 5489 milisegundos y el concurrente de 4399 milisegundos, lo que
representó una reducción media del 19,9\%. Las seis ejecuciones comparadas
completaron correctamente las tres fuentes. Estos resultados corresponden al
entorno local y al conjunto acotado de consultas del experimento; no deben
generalizarse como una garantía de rendimiento para cualquier carga o condición
de red.

No se implementó todavía paralelismo de CPU basado en múltiples procesos. Su
incorporación se considera una optimización futura para lotes masivos de
normalización o matching, siempre que una medición demuestre que esas tareas se
convierten en el cuello de botella.

\subsection{Implementación del Pipeline ETL}

El pipeline ETL se ejecuta después de cada extracción exitosa y mantiene una
traza auditable desde el dato crudo hasta el precio cargado. Sus etapas son:

\begin{enumerate}
    \item almacenamiento inmutable del resultado extraído en staging;
    \item validación de código externo, nombre y precio positivo;
    \item detección de códigos repetidos dentro de la misma corrida;
    \item limpieza de precios, nombres y presentaciones;
    \item normalización de unidades comparables, por ejemplo litros a
    mililitros y kilogramos a gramos;
    \item validación de GTIN-8, GTIN-12, GTIN-13 y GTIN-14 mediante su dígito
    verificador;
    \item resolución de identidad por publicación conocida, GTIN válido o
    equivalencia estructural no ambigua;
    \item creación controlada de productos y publicaciones faltantes;
    \item carga idempotente del historial de precios;
    \item registro del resultado como cargado, rechazado, duplicado o sin
    coincidencia.
\end{enumerate}

La identidad canónica conserva separada la entidad \textit{Producto}, utilizada
para comparar cadenas, de \textit{ProductoFuente}, que representa la publicación
original de cada supermercado. Los casos ambiguos no se fusionan automáticamente.
El sistema genera una cola de revisión asistida donde cada aprobación o rechazo
queda acompañado por evidencia, confianza, nota de decisión y fecha.

La normalización automática de categorías externas no está terminada. El
pipeline vigente normaliza nombres, precios, marcas cuando existe evidencia
consistente, unidades, presentaciones e identificadores; por lo tanto, no se
afirma todavía una unificación completa de categorías entre cadenas.

Antes de utilizar un precio en una comparación se aplica además una política de
calidad operativa. Por defecto se excluyen precios con más de catorce días y se
marcan como sospechosos los desvíos extremos respecto de la mediana histórica,
siempre que existan al menos dos observaciones anteriores. El ranking informa
si un producto quedó fuera por ausencia, antigüedad o anomalía.

\subsection{Implementación de la Base de Datos}

La demostración local utiliza SQLite y se encuentra versionada mediante siete
migraciones Alembic. El esquema incluye las entidades principales del dominio y
tablas adicionales para fuentes, corridas, staging e identidad:

\begin{itemize}
    \item provincias, ciudades, supermercados y sucursales;
    \item categorías, marcas, productos y productos por fuente;
    \item precios históricos;
    \item fuentes y corridas de scraping;
    \item productos extraídos en staging;
    \item revisiones auditables de identidad canónica.
\end{itemize}

SQLAlchemy abstrae la persistencia y la aplicación incluye controladores para
SQLite y PostgreSQL. PostgreSQL constituye el objetivo para producción debido a
su mejor comportamiento ante concurrencia, volumen y múltiples conexiones. Sin
embargo, la integración completa presentada y las pruebas automatizadas fueron
ejecutadas sobre SQLite; por ello, PostgreSQL debe considerarse compatible y
planificado, pero pendiente de una validación final equivalente.

\subsection{Implementación del Modelo Multicriterio}

El modelo de decisión utiliza el método de suma ponderada
(\textit{Weighted Sum Model}). Cada sucursal con cobertura completa de la
canasta se representa como una alternativa y se evalúa mediante costo total,
distancia y ahorro potencial.

La distancia se calcula con la fórmula de Haversine a partir de las coordenadas
de referencia y de la sucursal. El costo total considera la cantidad solicitada
de cada producto. El ahorro se obtiene como la diferencia entre el costo de la
alternativa más cara y el costo de cada sucursal.

Para una alternativa $i$, el ahorro potencial se define como:

\begin{equation}
Ahorro_i = Costo_{\max} - Costo_i
\label{eq:ahorro_potencial_alineado}
\end{equation}

Los criterios se normalizan antes de aplicar las ponderaciones. Para criterios
de costo, como precio y distancia, se utiliza:

\begin{equation}
N_{costo}(x_i) =
\frac{x_{\max} - x_i}{x_{\max} - x_{\min}}
\label{eq:normalizacion_costo_alineada}
\end{equation}

Para criterios de beneficio, como el ahorro, se utiliza:

\begin{equation}
N_{beneficio}(x_i) =
\frac{x_i - x_{\min}}{x_{\max} - x_{\min}}
\label{eq:normalizacion_beneficio_alineada}
\end{equation}

Cuando todos los valores de un criterio son iguales, el sistema asigna un valor
normalizado de uno a todas las alternativas para evitar una división por cero y
no introducir una diferencia inexistente.

El puntaje global se calcula mediante:

\begin{equation}
P_i = w_p N_{precio_i} + w_d N_{distancia_i} + w_a N_{ahorro_i}
\label{eq:puntaje_final_alineado}
\end{equation}

donde $w_p + w_d + w_a = 1$.

Los pesos iniciales son 0,6 para precio, 0,3 para distancia y 0,1 para ahorro,
aunque el usuario puede modificarlos desde la interfaz siempre que sean no
negativos y sumen uno. Las sucursales sin cobertura completa no reciben un
puntaje artificial: se excluyen del ranking principal y se muestran por separado
con los productos y motivos faltantes.

\subsection{Implementación del Frontend}

El frontend fue implementado con Vue.js 3 y TypeScript. La aplicación consume
la API REST en modo \textit{live} y conserva un modo de datos simulados para
desarrollo aislado.

La vista de comparación permite:

\begin{itemize}
    \item seleccionar una ciudad de referencia;
    \item buscar productos existentes en el catálogo;
    \item seleccionar las fuentes activas;
    \item ejecutar explícitamente una actualización en vivo para el texto
    buscado;
    \item agregar productos y cantidades a una canasta temporal;
    \item consultar precios vigentes y su antigüedad;
    \item configurar las ponderaciones del modelo;
    \item visualizar alternativas completas, puntajes, costos, distancias y
    ahorro;
    \item consultar sucursales excluidas y el motivo de la falta de cobertura.
\end{itemize}

La búsqueda común consulta primero el catálogo persistido. La extracción en
tiempo real no se dispara automáticamente ante cada pulsación: el usuario debe
utilizar la acción \textit{Actualizar precios}. Esta decisión evita solicitudes
innecesarias a sitios externos y permite seleccionar las fuentes involucradas.

La interfaz es adaptable a resoluciones de escritorio y dispositivos móviles.
Actualmente no incorpora un mapa ni solicita la geolocalización del navegador;
las distancias mostradas utilizan las coordenadas de la ciudad seleccionada como
origen. En consecuencia, la representación geográfica mediante mapas permanece
como una ampliación pendiente.

\subsection{Integración y Alcance Actual del Sistema}

El flujo integrado implementado comprende la configuración de fuentes, la
extracción concurrente, el staging, el ETL, la identidad canónica, el historial
de precios, la consulta desde la API, el cálculo del ranking y su visualización
en el frontend.

La solución debe caracterizarse como una demostración funcional avanzada para
una ciudad y un conjunto piloto de fuentes, y no como un servicio productivo de
cobertura nacional. En particular, permanecen pendientes:

\begin{itemize}
    \item validar o geocodificar todas las coordenadas de sucursales piloto;
    \item incorporar mapas y ubicación real del usuario;
    \item implementar la ejecución periódica mediante un scheduler operativo;
    \item validar migraciones, ETL y carga concurrente sobre PostgreSQL;
    \item ejecutar pruebas de carga, disponibilidad y despliegue;
    \item ampliar la evaluación experimental de calidad y rendimiento.
\end{itemize}

La actualización actual puede iniciarse desde la interfaz, la API o la línea de
comandos. Existe la estructura destinada al scheduler, pero el trabajo periódico
todavía no fue implementado. Por lo tanto, el sistema permite actualización bajo
demanda, no actualización automática programada.

Con este alcance, la implementación demuestra el recorrido técnico central de la
tesis y permite evaluar el modelo DSS con datos reales acotados, manteniendo
explícitas las limitaciones que deben resolverse antes de una puesta en producción.
